\icode{N22J-MCC-SI-L2-261023-001}
\nceeID[3]{2022北京}{8}
我国航天员在“天宫课堂”中演示了多种有趣的实验，提高了青少年科学探索的兴趣。某同学设计了如下实验：细绳一端固定，另一端系一小球，给小球一初速度使其在竖直平面内做圆周运动。无论在“天宫”还是在地面做此实验\paren[]

\begin{tikzpicture}[>=stealth,font=\scriptsize]
	\draw[dashed,cyan] (0,0) circle(1.5);
	\draw[magenta] (0,0)--(0,-1.5);
	\shade[ball color=gray] (0,-1.5) circle(.1);
	\draw[->] (0,-1.8)--node[below]{\(v_0\)}(-.6,-1.8);
	\fill (0,0) circle(.05);
\end{tikzpicture}
\begin{choices}
	\item 小球的速度大小均发生变化
	\item 小球的向心加速度大小均发生变化
	\item 细绳的拉力对小球均不做功
	\item 细绳的拉力大小均发生变化
\end{choices}
\begin{answer}
	C\qquad 提示：小球其实是匀速圆周运动
\end{answer}